\documentclass[
	fontsize=12pt,          % default font size 12pt
	paper=a4,               % DIN A4 page format
	numbers=noenddot,       % remove dots behind chapter numbers (e.g. 1.5 not 1.5.)
	listof=totoc,           % add list of figures, tables, etc. to ToC
	listof=entryprefix,     % add entry name to figures, tables, etc.
	listof=nochaptergap,    % no chapter gap for figures, tables, etc.
	bibliography=totoc,     % add bibliography to ToC but without a chapter number
	parskip=half            % half line spacing between paragraphs
	openany                 % chapters can start on any page
]{scrbook}
\usepackage{tcolorbox}

\include{../../for_latex/preamble}
\title{\Huge\textbf{Springender Ball}}
\author{Luis Staudt}
\date{}

% Code-Highlighting für Java
\definecolor{codegray}{rgb}{0.5,0.5,0.5}
\definecolor{codepurple}{rgb}{0.58,0,0.82}
\definecolor{backcolour}{rgb}{0.95,0.95,0.92}
\definecolor{keywordcolor}{rgb}{0.8,0.47,0.13}

\lstdefinestyle{java}{
	language=,
	basicstyle=\ttfamily\small,
	keywordstyle=\bfseries\color{blue!70!black},
	stringstyle=\color{red!60!black},
	commentstyle=\itshape\color{gray},
	numbers=left,
	numberstyle=\tiny\color{gray},
	numbersep=8pt,
	frame=single,
	framerule=0.4pt,
	rulecolor=\color{gray!50},
	breaklines=true,
	showstringspaces=false,
	tabsize=4,
	xleftmargin=1.5em,
	framexleftmargin=1.5em,
	aboveskip=0.8em,
	belowskip=0.6em
}
\lstset{style=java}

\newtcolorbox{hinweis}[1][Hinweis]{
	colback=blue!5,
	colframe=blue!40!black,
	fonttitle=\bfseries,
	title=#1,
	boxrule=0.5pt,
	arc=2pt
}

\newtcolorbox{beispiel}[1][Beispiel]{
	colback=green!5,
	colframe=green!40!black,
	fonttitle=\bfseries,
	title=#1,
	boxrule=0.5pt,
	arc=2pt
}

\begin{document}

% --- Titelblock ---
\begin{center}
	\rule{\textwidth}{0.8pt}\\[0.6em]
	{\LARGE\bfseries Intensivtutorium}\\[0.3em]
	{\large Programmierung}\\[0.6em]
	\begin{tabular}{r l @{\hspace{2cm}} r l}
		\textbf{Semester:}      & Sommersemester 2026                    & \textbf{Tutor:}         & Luis Staudt \\
		\textbf{Studiengänge:}  & MIB, OMB, PEB, MVB         & & \\
	\end{tabular}\\[0.6em]
	\rule{\textwidth}{0.8pt}
\end{center}
\vspace{1em}

% --- Seitenkopf und -fuß (ab Seite 2) ---
\pagestyle{fancy}
\fancyhf{}
\lhead{\small Intensivtutorium -- Grundlagen der Programmierung}
\rhead{\small SS 2026}
\cfoot{\thepage}
\renewcommand{\headrulewidth}{0.4pt}


% =============================================
% Übungsaufgabe -- Springender Ball
% =============================================

\section*{Übungsaufgabe -- Springender Ball}

Ein Ball wird aus der Höhe~$h_0$ fallen gelassen.
Bei jedem Aufprall verliert er Energie, sodass er nach dem $n$-ten Sprung nur noch die Höhe
\[
	h_n = h_0 \cdot k^{\,n}
\]
erreicht, wobei $k$ (mit $0 < k < 1$) der \textbf{Rückprallfaktor} ist.
Der Ball gilt als \glqq zur Ruhe gekommen\grqq{}, sobald seine Sprunghöhe kleiner als $\varepsilon = 0{,}01~\text{m}$ ist.

\begin{center}
	\begin{tikzpicture}[>=Stealth, scale=1]
		% Boden
		\draw[line width=0.8pt] (-0.3,0) -- (8.3,0);
		\foreach \x in {-0.1,0.2,...,8.2} {
			\draw[line width=0.3pt, gray] (\x,0) -- (\x-0.18,-0.22);
		}
		% Sprungbahn (abnehmende Parabelboegen)
		\draw[blue!60!black, line width=1pt] (0,3) .. controls (0.5,3) and (1.0,0) .. (1.5,0);
		\draw[blue!60!black, line width=1pt] (1.5,0) .. controls (2.1,2.1) and (2.7,2.1) .. (3.3,0);
		\draw[blue!60!black, line width=1pt] (3.3,0) .. controls (3.75,1.4) and (4.2,1.4) .. (4.65,0);
		\draw[blue!60!black, line width=1pt] (4.65,0) .. controls (4.95,0.85) and (5.25,0.85) .. (5.55,0);
		\draw[blue!60!black, line width=1pt] (5.55,0) .. controls (5.75,0.5) and (5.95,0.5) .. (6.15,0);
		% Ball
		\fill[red!70!black] (0,3) circle (0.12);
		% Hoehen
		\draw[<->, line width=0.5pt] (-0.15,0) -- (-0.15,3);
		\node[fill=white, inner sep=1pt] at (-0.15,1.5) {$h_0$};
		\draw[dotted] (2.4,0) -- (2.4,2.1);
		\draw[<->, line width=0.5pt] (2.4,0) -- (2.4,2.1);
		\node[fill=white, inner sep=1pt, font=\footnotesize] at (2.4,1.05) {$h_1$};
		\draw[dotted] (3.975,0) -- (3.975,1.4);
		\draw[<->, line width=0.5pt] (3.975,0) -- (3.975,1.4);
		\node[fill=white, inner sep=1pt, font=\footnotesize] at (3.975,0.7) {$h_2$};
	\end{tikzpicture}
\end{center}

Die auf dem gesamten Weg \textbf{zurückgelegte Gesamtstrecke} setzt sich zusammen aus dem ersten Fall (Strecke~$h_0$) und je Sprung der Auf- und der Abwärtsstrecke (also $2\,h_n$ pro Sprung).

\begin{hinweis}[Physikalischer Hintergrund]
	Der Rückprallfaktor~$k$ beschreibt, welcher Anteil der Höhe nach einem Aufprall erhalten bleibt.
	Ein kleines~$k$ bedeutet einen stark dämpfenden Ball (wenige, niedrige Sprünge), ein~$k$ nahe~1 einen sehr elastischen Ball (viele Sprünge).
	Die Höhen bilden eine geometrische Folge; die Gesamtstrecke bleibt trotz unendlich vieler theoretischer Sprünge endlich.
\end{hinweis}

\subsection*{Aufgabe}

Schreiben Sie ein Programm, das für \textbf{mehrere Fallhöhen} ($h_0 = 1, 2, 3~\text{m}$) und \textbf{mehrere Rückprallfaktoren} ($k = 0{,}5\,;\ 0{,}6\,;\ 0{,}7$) jeweils die folgenden beiden Größen bestimmt und für jede Kombination eine Zeile in die Datei \texttt{ball.txt} schreibt:

\begin{enumerate}[label=\textbf{\arabic*.)}]
	\item die \textbf{Anzahl der Sprünge}, bis der Ball zur Ruhe kommt (Sprunghöhe $< \varepsilon$),
	\item die insgesamt \textbf{zurückgelegte Gesamtstrecke}.
\end{enumerate}

Es sind also alle Kombinationen aus Fallhöhe und Rückprallfaktor zu durchlaufen (\textbf{zwei geschachtelte \texttt{for}-Schleifen}); die Simulation der Sprünge bis zum Stillstand erfolgt in einer \textbf{\texttt{while}-Schleife}.

\begin{hinweis}
	Legen Sie die beiden Wertelisten als Felder an, z.\,B.\ \lstinline|double[] fallhoehen = {1.0, 2.0, 3.0};|.
	Datei öffnen mit \lstinline|new FileWriter("ball.txt")|, je Zeile mit \texttt{write(\dots)} und \lstinline|"\n"| abschließen, am Ende \texttt{close()} nicht vergessen.
	Nutzen Sie \lstinline|"\t"| für die Spaltenabstände.
\end{hinweis}

\medskip
Vervollständige das folgende Programmgerüst:

\begin{lstlisting}
import java.io.FileWriter;
import java.io.IOException;

public class BouncingBall {

    public static void main(String[] args) throws IOException {

        double[] fallhoehen = {1.0, 2.0, 3.0};
        double[] rueckprallfaktoren = {0.5, 0.6, 0.7};
        double epsilon = 0.01;

        FileWriter writer = new FileWriter("ball.txt");
        // hier erweitern
    }
}
\end{lstlisting}

\begin{beispiel}[Erwarteter Dateiinhalt]
\begin{lstlisting}[language={}]
h0 [m]  k    Spruenge  Gesamtstrecke [m]
1.0	0.5	7	1.984375
1.0	0.6	10	2.9818601472
1.0	0.7	13	4.6214517951433995
2.0	0.5	8	3.984375
...
3.0	0.7	16	13.953473897202553
\end{lstlisting}
\end{beispiel}

\end{document}
